\subsection{Correlation heatmaps}
\label{sec:heatmaps}

Panels (a)--(c) of Fig.~\ref{fig:D04} show $S(\pi)$, the string correlation $O_{\rm str}$, and the combination
\begin{equation}
Q=\frac{4}{3}+2O_{\rm str}-\frac{1}{6}S(\pi).
\end{equation}
All three quantities are evaluated on the same parameter grid as the reflection marker, allowing comparison of the parameter dependence of staggered correlations, nonlocal string correlations, and their combination.

The general form of the string correlator is
\begin{equation}
O^z_{\rm str}(i,j)=\Big\langle S^z_i\exp\!\Big(i\pi\!\sum_{k=i+1}^{j-1}S^z_k\Big)S^z_j\Big\rangle,
\end{equation}
where $S^z_k=Z_k/2$. The spin-dependent phase factor probes hidden order within the interval, providing a diagnostic of valence-bond phases~\cite{dennijs1989preroughening,perezgarcia2008string}. For a finite chain, the measured quantity is the correlation between selected endpoints. On our $L=8$ chain we use $d=3$ with summed endpoint pairs, $O_{\rm str}=\langle(Z_1+Z_2)(Z_3Z_4)(Z_5Z_6)(Z_7+Z_8)\rangle$ (the two minus signs cancel for $L=8$); the heatmaps and the hardware protocol use $\hat S(\pi)=S(\pi)/8$ and $\hat O_{\rm str}=-O_{\rm str}$.

Both $S(\pi)$ and the $z$-string correlation can be evaluated from computational-basis bitstring statistics and serve as the main observables in our hardware protocol. The partial-reflection marker is used for the numerical phase diagram; measurement schemes for many-body invariants and randomized measurements are discussed in Refs.~\cite{elben2020manybody,elben2022toolbox}. The observable scales used in the hardware--simulator comparison are specified in Sec.~\ref{sec:comparison}.
